\section{zkSnark Design {\bf (NEW EXTENDED VERSION)}}
In this section, we discuss how to efficiently encode the
5-stage regex discharging algorithms in zkSNARK. Given that
there are many small files (e.g., emails), it is not economical
to provide a proof with size almost comparable to the original
file size. We opt for pairing based constant size proof system
such as Groth16 (instead of log-size proof system such as sumcheck
based systems). Our choice of proving schemes is also limited
by the choice of open-source folding scheme implementations available
at the time the project is implemented.
In this case, we choose to adopt Sonobe (an implementation of
Nova + Cyclefold folding schemes that work for BN254-Grumpkin curve
that allows producing final decider proof in Groth16). 
It is possible to encode our 5-stage discharging algorithms using
many other schemes, such as Plonk based for universal prover keys,
or Mangrove for PCD foldings schemes with constant memory footprints,
and we leave them for future work.

We have several design challenges to overcome.
First of all, we would like to take advantage of
the folding schemes so that the prover cost can be
minimized. Second, in addition to mass processing regex signatures,
we can also mass-discharge files, e.g., providing
a commitment to a large collection of files (e.g., all binary executables
or all Enron emails), provide an $O(1)$ proof for free of
malware for all of them, and then provide a membership 
of an individual file commitment belong to the combined commitment.
Third, to achieve the aforementioned scheme, we have to
handle vast difference in file sizes from $2^{9}$ to $2^{27}$.
This concerns handling non-uniform circuits (with different size).
Fourthly, we need to handle lookup arguments effectively.
Although CQ lookup is handled in ProtoStar, we present
a simpler stand-alone zk-CQ protocol which shows similar
concrete prover effiicency.

\subsection{Consideration of Folding Cost}
For the vast majority of folding schemes that follow the Nova folding,
it is shown that prover concrete efficiency is $4\times$ of Groth16.
However, there is a final decider scheme (in Sonobe) that is expensive
to compute, its cost is shown in the following, given $n$ the step-wise
circuit size, $k$ the number of steps applied in folding: 
\[
	2^{23.5} + 3n
\]
Mainly, the $2^{23.5}$ is a fixed non-native field arithmetic component in the decider and $3n$ is the linear cost on proving the step-wise R1CS relation.
Now given that folding itself takes $2n$ group operations per step, the fraction of the decider proof genreation among the total cost is, when $k$ is large enough:
\[
\frac{2^{23.5} + 3n}{2nk} \approx \frac{1.5}{k}
\]

In our later presentation, we will need to customize the Sonobe final decider
proof, which will add several more ``expensive" constant size computation. 
However, given that $k$ is large enough, they will be amortized soon, which
does not cause concern. 

Also another interesting point is the step-wise circuit size.
In the celebrative work of Mangrove, it is shown how any arbitrary Plonk
relation can be chunked into any arbitrary thread of smaller computation,
which allows arbitrary folds of scalability via parallelism, i.e.,
it allows chunking a large R1CS into many very small steps.
In our case, server RAM is not a concern, so that we allow large
step-wise circuit size, as long as it does not make $k$ too small.

\subsection{Road Map}
In this section, we will first discuss how to encode the $5$-stage 
regex discharging algorithm efficiently in R1CS. We then discuss a 
customized folding scheme to support batch process a large collection
files. 

\subsection{Encoding $5$-Stage Regex Discharging Algorithm}
\subsubsection{Encode DFA (TODO)}
In practice, an ACDFA can be transformed into a DFA. As DFA acceptance is
needed in various parts of the entire algorithm, we discuss them first.

The key idea is that: there is an external and precomputed lookup table
of all DFAs needed (notice that there are multiple tables). But we integrate
them into one table of two columns $\texttt{TableID}, \texttt{Entry}$.

TODOS in details: rough idea

(1) discuss how to encode DFA transition

(2) discuss what kind of information will be stored in external table, e.g., transitions (for different DFAs), the critical patterns needed by each signature).

(3) present an algorithm in R1CS or circuit, which takes in witnesses: (source state, char, dest state) sequence, produces (still secret) output: the set of signatures that should be triggered (those who have their critical patterns appear along the acceptance path)

(4) present a variety of gadgets that proves permutation of vectors, equivalence of vectors and its sorted form, packing a vector to a set etc. In particular, notice that we can use the in-circuit form of the Habock scheme \cite{Hab22} for proving lookup relation between two vectors (when they are small) {\em in circuit}, this is useful for proving a ``set" is reduced from a vector with duplicate items.

(5) Given the gadgets, this motivates the definition of R1CS relation with $2$ stages (to allow to generate the random needed by the gadget), and 
the need to tag some some elements for the use of external lookup argument.
